\section{System}
ContextTrace accepts a JSON trace containing query, answer, contexts, citations,
and optional metadata. The implementation normalizes the trace, decomposes the
answer into claims, discovers candidate evidence, evaluates support and
contradiction, checks citation alignment, assesses source metadata, decides
whether abstention was required, and infers a primary root cause. Each stage
returns typed JSON consumed by the next stage and by the final report.

Claim verification combines lexical anchors with deterministic semantic
normalization for entities, quantities, dates, negation, and common relational
paraphrases. Contradiction checks are independent of citation checks: a claim
may be contradicted by one source while citing a different source, or it may be
supported by the retrieved set while citing the wrong source. Evidence-span
localization records the source ID and minimal matched text. Source assessment
uses freshness, authority, and canonical-source metadata only when provided.

The diagnosis layer aggregates claim records into observable failure types and
a primary root cause. Repair planning ties a suggested change to cited evidence
and emits a regression-test target rather than claiming the change has already
worked. Production profiles expose citation, contradiction, abstention, source,
root-cause, and span modules as explicit switches; ablations therefore execute
the shipping verifier rather than a parallel research implementation.

The default pipeline is local, deterministic, and free of API cost. Model-based
judges are optional. A provider/model/prompt/cache/cost tuple must be frozen
before such outputs enter a result table. No qualifying NLI-only or judge-only
configuration existed at freeze time, so both remain registered as unavailable.
